% УВОД, без номер.
\section*{УВОД}
\label{sec:introduction}
\addcontentsline{toc}{section}{УВОД}

Наблюдението на работеща система е задача, при която закъснението има цена. Ако
показателят стигне до оператора със закъснение от няколко секунди, картината на екрана
описва състояние, което вече не съществува. Класическият уеб модел, при който браузърът
пита сървъра през равни интервали и получава пълна страница в отговор, харчи заявки,
когато нищо не се е променило, и въпреки това изостава, когато промяната настъпи между
две запитвания. Двупосочната връзка, при която сървърът изпраща стойност в момента, в
който я узнае, решава и двата проблема наведнъж.

Настоящият курсов проект разработва \textbf{Pulse}: уеб система за наблюдение на
показатели в реално време. Сървърната част е програма на C++20, която сама обслужва
HTTP/1.1~\cite{rfc9112} за първоначалното зареждане и WebSocket~\cite{rfc6455} за потока
от стойности след това. Клиентската част е страница без приложна рамка, която стъпва само
на браузърните стандарти: обектния модел на документа~\cite{dom_ls}, модела на събитията,
структурната графика във формат SVG~\cite{svg2} и разположението с CSS
Grid~\cite{css_grid1}.

\textbf{Целта} е работеща система за наблюдение в реално време, при която всяка тема от
учебното съдържание на дисциплината е представена от действащ елемент на системата, а не
от отделен демонстрационен пример. Задачите са: анализ на технологиите за обмен между
браузър и сървър и обосновка на избора между тях; проектиране на архитектурата и на
договора между двете страни; програмна реализация на сървъра на C++20 и на клиента без
рамка; и измерване на закъснението, на обема на съобщението и на цената на двата формата
за обмен, с анализ на получените числа.

\textbf{Обхватът} е един процес на сървъра, който събира показатели и ги раздава, и една
страница в браузъра, която ги изобразява. Извън обхвата остават разпределеното съхранение
на исторически редове, удостоверяването на потребители и разгръщането в производствена
среда.
